\chapter{Chapter08 — Product LTF Noise Correction}
%%% generated by the blueprint pipeline from FABL.Chapter08.ProductLTFNoiseCorrection
\section{Overview}

Declarations of the Lean module \texttt{FABL.Chapter08.ProductLTFNoiseCorrection}.

\section{Declarations}

\begin{definition}[Selector retention]
\lean{FABL.selectorRetention}
\label{FABL.selectorRetention}
\leanok
For a real number $\theta$, define the selector retention to be $(1+\theta)/2$.
\end{definition}

\begin{theorem}[Selector retention probability bounds]
\lean{FABL.selectorRetention_mem_Icc}
\label{FABL.selectorRetention_mem_Icc}
\leanok
\uses{FABL.selectorRetention}
\difficulty{2}
If $\theta\in[-1,1]$, then $(1+\theta)/2\in[0,1]$.
\end{theorem}

\begin{definition}[Retention correlation]
\lean{FABL.retentionCorrelation}
\statementsource{Analysis-of-Boolean-Functions-by-Ryan-ODonnell}{pdf-ffbe9c7b3306-p214-b005}
\label{FABL.retentionCorrelation}
\leanok
For a real number $\rho$, its retention correlation is $2\rho-1$.
\end{definition}

\begin{theorem}[Range of the retention correlation]
\lean{FABL.retentionCorrelation_mem_Icc}
\label{FABL.retentionCorrelation_mem_Icc}
\leanok
\uses{FABL.retentionCorrelation}
\difficulty{2}
For every retention parameter $\rho\in[0,1]$, its retention correlation $2\rho-1$ lies
in $[-1,1]$.
\end{theorem}

\begin{theorem}[Retention parameter inversion]
\lean{FABL.selectorRetention_retentionCorrelation}
\label{FABL.selectorRetention_retentionCorrelation}
\leanok
\uses{FABL.retentionCorrelation, FABL.selectorRetention}
\difficulty{3}
For every real number $\rho$, applying the selector-retention function to the retention
correlation $2\rho-1$ returns $\rho$.
\end{theorem}

\begin{theorem}[Inverse retention parametrization]
\lean{FABL.retentionCorrelation_selectorRetention}
\label{FABL.retentionCorrelation_selectorRetention}
\leanok
\uses{FABL.retentionCorrelation, FABL.selectorRetention}
\difficulty{3}
For every real parameter $\theta$, applying the retention correlation $2\rho-1$ to the
retention probability $(1+\theta)/2$ gives $\theta$.
\end{theorem}

\begin{definition}[Selector product value]
\lean{FABL.selectProductValue}
\label{FABL.selectProductValue}
\leanok
\uses{FABL.Sign}
Given a type $\Omega$, two elements $a,b\in\Omega$, and a sign $s\in\{\pm1\}$, the
selector product value is $a$ when $s=1$ and $b$ when $s=-1$.
\end{definition}

\begin{theorem}[Fourier expansion of a function on two signs]
\lean{FABL.signPair_fourier_expansion}
\label{FABL.signPair_fourier_expansion}
\leanok
\uses{FABL.Sign, FABL.signValue}
\difficulty{4}
Let $H:\{\pm1\}^2\to\mathbb R$, and write $z=(z_1,z_2)\in\{\pm1\}^2$. Then
\[
H(z)=\frac{H(1,1)+H(1,-1)+H(-1,1)+H(-1,-1)}4+\frac{H(1,1)+H(1,-1)-H(-1,1)-H(-1,-1)}4z_1+\frac{H(1,1)-H(1,-1)+H(-1,1)-H(-1,-1)}4z_2+\frac{H(1,1)-H(1,-1)-H(-1,1)+H(-1,-1)}4z_1z_2.
\]
\end{theorem}

\begin{theorem}[Expectation of a correlated sign pair]
\lean{FABL.pmfExpectation_correlatedSignPairPMF_eq_four_values}
\label{FABL.pmfExpectation_correlatedSignPairPMF_eq_four_values}
\leanok
\uses{FABL.Sign, FABL.correlatedSignPairPMF, FABL.correlatedSignPairPMF_expect_signValue_fst, FABL.correlatedSignPairPMF_expect_signValue_mul, FABL.correlatedSignPairPMF_expect_signValue_snd, FABL.pmfExpectation, FABL.pmfExpectation_add_fun, FABL.pmfExpectation_const, FABL.pmfExpectation_const_mul, FABL.signPair_fourier_expansion, FABL.signValue}
\difficulty{5}
Let $\theta\in[-1,1]$, and let $H$ be a real-valued function on pairs of signs. Under
the distribution that chooses a uniform sign $X$ and then sets $Y=X$ with probability
$(1+\theta)/2$ and $Y=-X$ otherwise, the expectation of $H(X,Y)$ is
$\bigl((1+\theta)(H(1,1)+H(-1,-1))+(1-\theta)(H(1,-1)+H(-1,1))\bigr)/4$.
\end{theorem}

\begin{definition}[Selector coordinate pair distribution]
\lean{FABL.selectorCoordinatePairPMF}
\label{FABL.selectorCoordinatePairPMF}
\leanok
\uses{FABL.correlatedSignPairPMF, FABL.selectProductValue}
Let $\Omega$ be a finite set, let $\pi$ be a probability distribution on $\Omega$, and
let $\theta\in[-1,1]$. Draw two independent elements from $\pi$, then draw a correlated
pair of signs with parameter $\theta$; apply each sign to the two drawn elements by
selecting the first element for sign $1$ and the second for sign $-1$. The resulting
pair of elements of $\Omega$ has the selector coordinate pair distribution.
\end{definition}

\begin{definition}[Coordinate retention pair distribution]
\lean{FABL.productRetentionCoordinatePairPMF}
\statementsource{Analysis-of-Boolean-Functions-by-Ryan-ODonnell}{pdf-ffbe9c7b3306-p214-b005}
\label{FABL.productRetentionCoordinatePairPMF}
\leanok
\uses{FABL.productCoordinateNoisePMF}
Let $\Omega$ be a finite type, let $\pi$ be a probability mass function on $\Omega$, and
let $\rho\in[0,1]$. The coordinate retention pair distribution on $\Omega\times\Omega$
is obtained by drawing $x$ according to $\pi$, then drawing $y$ from the coordinate
product noise distribution with base law $\pi$, retention parameter $\rho$, and input
$x$, and returning $(x,y)$.
\end{definition}

\begin{theorem}[Linearity of expectation]
\lean{FABL.pmfExpectation_two_linear}
\label{FABL.pmfExpectation_two_linear}
\leanok
\uses{FABL.pmfExpectation, FABL.pmfExpectation_add_fun, FABL.pmfExpectation_const_mul}
\difficulty{1}
Let $\Omega$ be a finite set, let $\pi$ be a probability mass function on $\Omega$, let
$a,b\in\mathbb{R}$, and let $f,g:\Omega\to\mathbb{R}$. Then
$\mathbb{E}_{\pi}[af+bg]=a\mathbb{E}_{\pi}[f]+b\mathbb{E}_{\pi}[g]$.
\end{theorem}

\begin{theorem}[Linearity of finite probability expectation]
\lean{FABL.pmfExpectation_four_linear}
\label{FABL.pmfExpectation_four_linear}
\leanok
\uses{FABL.pmfExpectation, FABL.pmfExpectation_add_fun, FABL.pmfExpectation_const_mul}
\difficulty{3}
For a finite set $\Omega$, a probability mass function $\pi$ on $\Omega$, real scalars
$a,b,c,d$, and functions $f,g,h,k:\Omega\to\mathbb{R}$, the expectation of their linear
combination is
$\mathbb{E}_{\pi}[af+bg+ch+dk]=a\mathbb{E}_{\pi}[f]+b\mathbb{E}_{\pi}[g]+c\mathbb{E}_{\pi}[h]+d\mathbb{E}_{\pi}[k]$.
\end{theorem}

\begin{theorem}[Interchanging two finite expectations]
\lean{FABL.pmfExpectation_swap}
\label{FABL.pmfExpectation_swap}
\leanok
\uses{FABL.pmfExpectation}
\difficulty{4}
Let $\Omega$ be a finite set, let $\pi$ be a probability mass function on $\Omega$, and
let $H:\Omega\times\Omega\to\mathbb{R}$. The two iterated expectations agree when the
order of averaging is reversed:
$\mathbb{E}_{x\sim\pi}\mathbb{E}_{y\sim\pi}[H(x,y)]=\mathbb{E}_{y\sim\pi}\mathbb{E}_{x\sim\pi}[H(x,y)]$.
\end{theorem}

\begin{theorem}[Equality of two selector pair distributions]
\lean{FABL.pmfExpectation_selectorCoordinatePairPMF}
\label{FABL.pmfExpectation_selectorCoordinatePairPMF}
\leanok
\uses{FABL.pmfExpectation, FABL.pmfExpectation_bind, FABL.pmfExpectation_const, FABL.pmfExpectation_correlatedSignPairPMF_eq_four_values, FABL.pmfExpectation_four_linear, FABL.pmfExpectation_map, FABL.pmfExpectation_productCoordinateNoisePMF, FABL.pmfExpectation_swap, FABL.pmfExpectation_two_linear, FABL.productRetentionCoordinatePairPMF, FABL.selectProductValue, FABL.selectorCoordinatePairPMF, FABL.selectorRetention, FABL.selectorRetention_mem_Icc}
\difficulty{6}
Let $\Omega$ be a finite set, let $\pi$ be a probability mass function on $\Omega$, let
$\theta\in[-1,1]$, and let $H:\Omega\times\Omega\to\mathbb{R}$. Draw two independent
values from $\pi$ and two signs in $\{\pm1\}$ so that the first sign is uniform and the
second equals it with probability $(1+\theta)/2$ and is its negative otherwise; use each
sign to select the first or second drawn value, respectively, and apply $H$ to the
resulting pair. The expected value equals that obtained by drawing $x$ from $\pi$, then
retaining $x$ as the second value with probability $(1+\theta)/2$ and otherwise drawing
that second value independently from $\pi$, and applying $H$ to the pair.
\end{theorem}

\begin{theorem}[Selector law as product retention]
\lean{FABL.selectorCoordinatePairPMF_eq_productRetention}
\label{FABL.selectorCoordinatePairPMF_eq_productRetention}
\leanok
\uses{FABL.pmfExpectation, FABL.pmfExpectation_selectorCoordinatePairPMF, FABL.productRetentionCoordinatePairPMF, FABL.selectorCoordinatePairPMF, FABL.selectorRetention, FABL.selectorRetention_mem_Icc}
\difficulty{4}
Let $\Omega$ be a finite type, let $\pi$ be a probability mass function on $\Omega$, and
let $\theta\in[-1,1]$. The joint distribution formed by drawing two independent values
from $\pi$ and selecting between them using a correlated sign pair with parameter
$\theta$ equals the joint distribution formed by drawing the first value from $\pi$ and,
for the second coordinate, retaining it with probability $(1+\theta)/2$ and otherwise
drawing independently from $\pi$.
\end{theorem}

\begin{definition}[Product pair coordinate equivalence]
\lean{FABL.productPairCoordinatesEquiv}
\label{FABL.productPairCoordinatesEquiv}
\leanok
For a type $\Omega$ and a natural number $n$, the pairs of functions
$u,v:\{0,\ldots,n-1\}\to\Omega$ are in bijection with functions
$z:\{0,\ldots,n-1\}\to\Omega\times\Omega$, by sending $(u,v)$ to the function
$i\mapsto(u(i),v(i))$.
\end{definition}

\begin{definition}[Product selector pair distribution]
\lean{FABL.productSelectorPairPMF}
\label{FABL.productSelectorPairPMF}
\leanok
\uses{FABL.productPairCoordinatesEquiv, FABL.productProbabilityPMF, FABL.selectorCoordinatePairPMF}
Let $\Omega$ be finite, let $\pi$ be a probability distribution on $\Omega$, let $n$ be
a natural number, and let $\theta\in[-1,1]$. Draw $n$ independent coordinate pairs from
the selector coordinate pair distribution with base distribution $\pi$ and parameter
$\theta$, then regard these pairs as two $n$-tuples in $\Omega$. Their joint
distribution is the product selector pair distribution.
\end{definition}

\begin{definition}[Product retention pair distribution]
\lean{FABL.productRetentionPairPMF}
\statementsource{Analysis-of-Boolean-Functions-by-Ryan-ODonnell}{pdf-ffbe9c7b3306-p214-b005}
\label{FABL.productRetentionPairPMF}
\leanok
\uses{FABL.productNoiseKernel, FABL.productProbabilityPMF}
Let $\Omega$ be a finite type, let $\pi$ be a probability mass function on $\Omega$, let
$n$ be a natural number, and let $\rho\in[0,1]$. The product retention pair distribution
on $\Omega^n\times\Omega^n$ is obtained by drawing $x$ from the product probability
distribution with base law $\pi$ on $n$ coordinates, then drawing $y$ from the product
noise kernel with base law $\pi$, retention parameter $\rho$, and input $x$, and
returning $(x,y)$.
\end{definition}

\begin{theorem}[Joint coordinate-pair mass]
\lean{FABL.productRetentionCoordinatePairPMF_apply}
\label{FABL.productRetentionCoordinatePairPMF_apply}
\leanok
\uses{FABL.productCoordinateNoisePMF, FABL.productRetentionCoordinatePairPMF}
\difficulty{4}
Let $\Omega$ be a finite set, let $\pi$ be a probability mass function on $\Omega$, and
let $\rho\in[0,1]$. For any $x,y\in\Omega$, the joint probability mass at $(x,y)$ under
the retention-and-resampling experiment is $\pi(x)$ times the conditional probability
mass of obtaining $y$ from $x$ under the coordinate noise distribution.
\end{theorem}

\begin{theorem}[Joint product-pair mass]
\lean{FABL.productRetentionPairPMF_apply}
\label{FABL.productRetentionPairPMF_apply}
\leanok
\uses{FABL.productNoiseKernel, FABL.productProbabilityPMF, FABL.productRetentionPairPMF}
\difficulty{4}
Let $\Omega$ be a finite set, let $\pi$ be a probability mass function on $\Omega$, let
$n\ge 0$, and let $\rho\in[0,1]$. For sequences $x,y\in\Omega^n$, their joint mass under
the retention-and-resampling experiment is the mass of $x$ under the product
distribution $\pi^{\otimes n}$ multiplied by the conditional mass of $y$ given $x$ under
the product noise kernel.
\end{theorem}

\begin{theorem}[Coordinatewise product representation]
\lean{FABL.productRetentionPairPMF_eq_coordinateProduct}
\label{FABL.productRetentionPairPMF_eq_coordinateProduct}
\leanok
\uses{FABL.productNoiseKernel, FABL.productPairCoordinatesEquiv, FABL.productProbabilityPMF, FABL.productProbabilityPMF_apply, FABL.productRetentionCoordinatePairPMF, FABL.productRetentionCoordinatePairPMF_apply, FABL.productRetentionPairPMF, FABL.productRetentionPairPMF_apply}
\difficulty{4}
Let $\Omega$ be a finite set, let $\pi$ be a probability mass function on $\Omega$, let
$n\ge 0$, and let $\rho\in[0,1]$. The joint law of two sequences generated by first
sampling the initial sequence from $\pi^{\otimes n}$ and then applying independent
coordinate noise with retention parameter $\rho$ equals the law obtained by
independently sampling $n$ coordinate pairs from the one-coordinate
retention-and-resampling joint distribution and arranging the first and second entries
into their respective sequences.
\end{theorem}

\begin{theorem}[Selector and retention pair laws]
\lean{FABL.productSelectorPairPMF_eq_productRetention}
\label{FABL.productSelectorPairPMF_eq_productRetention}
\leanok
\uses{FABL.productRetentionPairPMF, FABL.productRetentionPairPMF_eq_coordinateProduct, FABL.productSelectorPairPMF, FABL.selectorCoordinatePairPMF_eq_productRetention, FABL.selectorRetention, FABL.selectorRetention_mem_Icc}
\difficulty{3}
Let $\Omega$ be a finite set, let $\pi$ be a probability mass function on $\Omega$, let
$n\ge 0$, and let $\theta\in[-1,1]$. The law of two $\Omega^n$-valued sequences obtained
by independently sampling coordinatewise value pairs from $\pi$ and selecting between
each pair using correlated uniform signs of correlation $\theta$ equals the
retention-and-resampling pair law with retention parameter $(1+\theta)/2$.
\end{theorem}

\begin{theorem}[Retention parameter and selector correlation]
\lean{FABL.productSelectorPairPMF_retentionCorrelation}
\label{FABL.productSelectorPairPMF_retentionCorrelation}
\leanok
\uses{FABL.productRetentionPairPMF, FABL.productSelectorPairPMF, FABL.productSelectorPairPMF_eq_productRetention, FABL.retentionCorrelation, FABL.retentionCorrelation_mem_Icc, FABL.selectorRetention_retentionCorrelation}
\difficulty{1}
Let $\Omega$ be a finite set, let $\pi$ be a probability mass function on $\Omega$, let
$n\ge 0$, and let $\rho\in[0,1]$. The selector-based pair law on $\Omega^n$, with sign
correlation $2\rho-1$, equals the retention-and-resampling pair law with retention
parameter $\rho$.
\end{theorem}

\begin{theorem}[Binding an independent product coordinatewise]
\lean{FABL.independentProductPMF_bind_coordinates}
\label{FABL.independentProductPMF_bind_coordinates}
\leanok
\uses{FABL.independentProductPMF, FABL.independentProductPMF_apply}
\difficulty{3}
Let $I$ be a finite index set, and let $A_i$ and $B_i$ be finite sets for each $i\in I$.
Choose a probability mass function $p_i$ on $A_i$ and, for each $a\in A_i$, a
probability mass function $q_i(a)$ on $B_i$. If a tuple is first sampled from the
independent product of the $p_i$ and then each coordinate $i$ is sampled from $q_i$ at
its original value, the resulting distribution is the independent product of the
coordinatewise distributions obtained by sampling from $p_i$ and then from $q_i$.
\end{theorem}

\begin{theorem}[Mapping an independent product coordinatewise]
\lean{FABL.independentProductPMF_map_coordinates}
\label{FABL.independentProductPMF_map_coordinates}
\leanok
\uses{FABL.independentProductPMF, FABL.independentProductPMF_apply}
\difficulty{5}
Let $I$ be a finite index set, and let $A_i$ and $B_i$ be finite sets for each $i\in I$.
Choose a probability mass function $p_i$ on each $A_i$ and a function $f_i:A_i\to B_i$
for each coordinate. Mapping each coordinate of the independent product through its
corresponding $f_i$ gives the independent product of the distributions obtained by
mapping each $p_i$ through $f_i$.
\end{theorem}

\begin{theorem}[Binding product distributions coordinatewise]
\lean{FABL.productProbabilityPMF_bind_coordinates}
\label{FABL.productProbabilityPMF_bind_coordinates}
\leanok
\uses{FABL.independentProductPMF, FABL.independentProductPMF_bind_coordinates, FABL.productProbabilityPMF}
\difficulty{1}
Let $A$ and $B$ be finite sets, let $p$ be a probability mass function on $A$, and for
each $a\in A$ let $q(a)$ be a probability mass function on $B$. For every natural number
$n$, first sample an independent $n$-tuple from $p$ and then, independently at each
coordinate conditional on that tuple, sample from the corresponding law $q(a)$. The
resulting distribution on $B^n$ is the independent product distribution whose coordinate
law is the mixture of $q(a)$ over $a$ drawn from $p$.
\end{theorem}

\begin{definition}[Literal product selector coordinate distribution]
\lean{FABL.literalProductSelectorCoordinatePMF}
\label{FABL.literalProductSelectorCoordinatePMF}
\leanok
\uses{FABL.correlatedSignPairPMF, FABL.productProbabilityPMF, FABL.selectProductValue}
Let $\Omega$ be finite, let $\pi$ be a probability distribution on $\Omega$, let $n$ be
a natural number, and let $\theta\in[-1,1]$. Independently draw two $n$-tuples from the
product distribution generated by $\pi$, and draw an $n$-tuple of correlated sign pairs
with parameter $\theta$. At each coordinate, use the first sign to select between the
corresponding entries of the two sampled tuples, and use the second sign to make another
such selection. The resulting $n$-tuple of pairs has the literal product selector
coordinate distribution.
\end{definition}

\begin{definition}[Selector pair given a fixed coordinate value]
\lean{FABL.selectorGivenPlusCoordinatePMF}
\label{FABL.selectorGivenPlusCoordinatePMF}
\leanok
\uses{FABL.correlatedSignPairPMF, FABL.selectProductValue}
Let $Ω$ be a finite set, let $π$ be a probability mass function on $Ω$, and let
$θ	ext{ satisfy }-1	ext{le}θ	ext{le}1$. Fix $plus	ext{in}Ω$. Sample $minus$
according to $π$, then sample a pair of selectors from the $θ$-correlated sign-pair
law; applying each selector to the pair $(plus,minus)$ gives a probability mass
function on $Ω	imesΩ$.
\end{definition}

\begin{theorem}[Conditional selector law factors by coordinate]
\lean{FABL.literalSelector_after_plus}
\label{FABL.literalSelector_after_plus}
\leanok
\uses{FABL.correlatedSignPairPMF, FABL.independentProductPMF, FABL.independentProductPMF_bind_coordinates, FABL.independentProductPMF_map_coordinates, FABL.productProbabilityPMF, FABL.selectProductValue, FABL.selectorGivenPlusCoordinatePMF}
\difficulty{4}
Let $\Omega$ be a finite set, let $\pi$ be a probability distribution on $\Omega$, let
$n$ be a natural number, let $\theta\in[-1,1]$, and fix a vector $a\in\Omega^n$. Sample
a second vector independently from $\pi^{\otimes n}$ and, independently at each
coordinate, a sign pair whose first sign is uniform and whose second sign agrees with it
with probability $(1+\theta)/2$. Use each sign to select between the fixed coordinate of
$a$ and the corresponding coordinate of the sampled vector. The resulting law on pairs
of output vectors equals the product over coordinates of the one-coordinate law obtained
by this procedure with the corresponding fixed value $a_i$.
\end{theorem}

\begin{theorem}[Coordinatewise product law for selector pairs]
\lean{FABL.literalProductSelectorCoordinatePMF_eq_product}
\label{FABL.literalProductSelectorCoordinatePMF_eq_product}
\leanok
\uses{FABL.correlatedSignPairPMF, FABL.independentProductPMF, FABL.literalProductSelectorCoordinatePMF, FABL.literalSelector_after_plus, FABL.productProbabilityPMF, FABL.productProbabilityPMF_bind_coordinates, FABL.selectProductValue, FABL.selectorCoordinatePairPMF, FABL.selectorGivenPlusCoordinatePMF}
\difficulty{3}
Let $\Omega$ be a finite set, let $\pi$ be a probability distribution on $\Omega$, let
$n$ be a natural number, and let $\theta\in[-1,1]$. Sample two independent vectors with
law $\pi^{\otimes n}$ and, independently at each coordinate, a pair of signs whose first
sign is uniform and whose second sign equals the first with probability $(1+\theta)/2$.
At each coordinate, use the first sign to select between the two sampled values for the
first output and the second sign for the second output. The resulting law on vectors of
value pairs equals the product, over coordinates, of the one-coordinate law obtained by
this same procedure.
\end{theorem}

\begin{definition}[Literal product selector pair law]
\lean{FABL.literalProductSelectorPairPMF}
\label{FABL.literalProductSelectorPairPMF}
\leanok
\uses{FABL.literalProductSelectorCoordinatePMF, FABL.productPairCoordinatesEquiv}
Let $Ω$ be a finite set, let $π$ be a probability mass function on $Ω$, let $n$ be a
natural number, and let $θ\in[-1,1]$. The literal product selector coordinate law,
transported by the coordinate-pair equivalence, is a probability mass function on
$(Ω^n)\times(Ω^n)$.
\end{definition}

\begin{theorem}[Two constructions of selector pair laws]
\lean{FABL.literalProductSelectorPairPMF_eq_productSelectorPairPMF}
\label{FABL.literalProductSelectorPairPMF_eq_productSelectorPairPMF}
\leanok
\uses{FABL.literalProductSelectorCoordinatePMF_eq_product, FABL.literalProductSelectorPairPMF, FABL.productSelectorPairPMF}
\difficulty{1}
Let $\Omega$ be a finite set, let $\pi$ be a probability distribution on $\Omega$, let
$n$ be a natural number, and let $\theta\in[-1,1]$. The law formed by sampling two
independent vectors from $\pi^{\otimes n}$ and applying independent coordinatewise
correlated sign pairs to select the two outputs is the same as the law formed by
independently sampling, at each coordinate, the one-coordinate pair distribution
obtained by that procedure.
\end{theorem}

\begin{theorem}[Selector pairs and product retention]
\lean{FABL.literalProductSelectorPairPMF_eq_productRetention}
\label{FABL.literalProductSelectorPairPMF_eq_productRetention}
\leanok
\uses{FABL.literalProductSelectorPairPMF, FABL.literalProductSelectorPairPMF_eq_productSelectorPairPMF, FABL.productRetentionPairPMF, FABL.productSelectorPairPMF_eq_productRetention, FABL.selectorRetention, FABL.selectorRetention_mem_Icc}
\difficulty{1}
Let $\Omega$ be a finite set, let $\pi$ be a probability distribution on $\Omega$, let
$n$ be a natural number, and let $\theta\in[-1,1]$. Sample $x$ from $\pi^{\otimes n}$,
then independently at each coordinate retain $x_i$ with probability $(1+\theta)/2$ and
otherwise draw a fresh value from $\pi$. The law of the pair of vectors $(x,y)$ equals
the law obtained by sampling two independent vectors from $\pi^{\otimes n}$ and
applying, coordinatewise, a correlated pair of signs with uniform first sign and second
sign agreeing with it with probability $(1+\theta)/2$ to select between the two values.
\end{theorem}

\begin{theorem}[Expectation under product selector sampling]
\lean{FABL.pmfExpectation_literalProductSelectorPairPMF}
\label{FABL.pmfExpectation_literalProductSelectorPairPMF}
\leanok
\uses{FABL.correlatedPairPMF, FABL.correlatedPairPMF_map_pairCoordinatesEquiv, FABL.correlatedSignPairPMF, FABL.literalProductSelectorCoordinatePMF, FABL.literalProductSelectorPairPMF, FABL.pairCoordinatesEquiv, FABL.pmfExpectation, FABL.pmfExpectation_bind, FABL.pmfExpectation_map, FABL.productProbabilityPMF, FABL.selectProductInput}
\difficulty{4}
Let $\Omega$ be a finite set with probability mass function $\pi$, let $n\geq 0$, and
let $\theta\in[-1,1]$. Independently sample two vectors $U,V\in\Omega^n$ by drawing each
coordinate from $\pi$. Also sample a pair of sign vectors
$(R,T)\in\{-1,1\}^n\times\{-1,1\}^n$ from the correlated sign-pair distribution with
parameter $\theta$. Form vectors $X,Y\in\Omega^n$ coordinatewise by setting $X_i=U_i$
when $R_i=1$ and $X_i=V_i$ otherwise, and setting $Y_i=U_i$ when $T_i=1$ and $Y_i=V_i$
otherwise. For every $H:\Omega^n\times\Omega^n\to\mathbb{R}$, its expectation at the
pair $(X,Y)$ equals the nested expectation obtained by first averaging over $U$, then
over $V$, and then over $(R,T)$.
\end{theorem}

\begin{theorem}[Complement of a unit interval parameter]
\lean{FABL.one_sub_mem_Icc}
\label{FABL.one_sub_mem_Icc}
\leanok
\difficulty{2}
If $\delta$ is a real number with $0\leq\delta\leq 1$, then $1-\delta$ also lies in
$[0,1]$.
\end{theorem}

\begin{definition}[Product resampling noise sensitivity]
\lean{FABL.productResamplingNoiseSensitivity}
\statementsource{Analysis-of-Boolean-Functions-by-Ryan-ODonnell}{
  pdf-ffbe9c7b3306-p053-b003,
  pdf-ffbe9c7b3306-p054-b002,
  pdf-ffbe9c7b3306-p130-b004,
  pdf-ffbe9c7b3306-p214-b005,
  pdf-ffbe9c7b3306-p241-b003
}
\label{FABL.productResamplingNoiseSensitivity}
\leanok
\uses{FABL.Sign, FABL.one_sub_mem_Icc, FABL.pmfExpectation, FABL.productRetentionPairPMF}
Let $\Omega$ be a finite type, let $\pi$ be a probability mass function on $\Omega$, let
$n$ be a natural number, and let $\delta\in[0,1]$. For a function
$f:\Omega^n\to\{\pm1\}$, its product resampling noise sensitivity is the probability,
under the product retention pair distribution with retention parameter $1-\delta$, that
the two sampled inputs receive different values under $f$.
\end{definition}

\begin{theorem}[Product resampling noise sensitivity le peres]
\lean{FABL.productResamplingNoiseSensitivity_le_peres}
\statementsource{Analysis-of-Boolean-Functions-by-Ryan-ODonnell}{
  pdf-ffbe9c7b3306-p130-b004,
  pdf-ffbe9c7b3306-p241-b003
}
\label{FABL.productResamplingNoiseSensitivity_le_peres}
\leanok
\uses{FABL.BooleanFunction, FABL.IsLinearThreshold, FABL.IsProductLinearThreshold, FABL.Sign, FABL.correlatedPairPMF, FABL.isLinearThreshold_selectProductInput, FABL.literalProductSelectorPairPMF, FABL.literalProductSelectorPairPMF_eq_productRetention, FABL.noiseSensitivity, FABL.one_sub_mem_Icc, FABL.one_sub_two_mul_mem_Icc, FABL.peresNoiseSensitivityBound, FABL.pmfExpectation, FABL.pmfExpectation_const, FABL.pmfExpectation_literalProductSelectorPairPMF, FABL.pmfExpectation_mono, FABL.productProbabilityPMF, FABL.productResamplingNoiseSensitivity, FABL.productRetentionPairPMF, FABL.selectProductInput, FABL.selectorRetention}
\difficulty{5}
% dataset-pipeline: stub -- prose pending build_informal_statements.py
Statement pending regeneration.
\end{theorem}

\begin{theorem}[Zero resampling noise sensitivity]
\lean{FABL.productResamplingNoiseSensitivity_zero}
\label{FABL.productResamplingNoiseSensitivity_zero}
\leanok
\uses{FABL.BooleanFunction, FABL.Sign, FABL.correlatedPairPMF, FABL.literalProductSelectorPairPMF, FABL.literalProductSelectorPairPMF_eq_productRetention, FABL.noiseSensitivity, FABL.noiseSensitivity_eq_sum_level, FABL.pmfExpectation, FABL.pmfExpectation_const, FABL.pmfExpectation_literalProductSelectorPairPMF, FABL.productProbabilityPMF, FABL.productResamplingNoiseSensitivity, FABL.productRetentionPairPMF, FABL.selectProductInput, FABL.selectorRetention}
\difficulty{5}
Let $\Omega$ be a finite set, let $\pi$ be a probability mass function on $\Omega$, let
$n$ be a natural number, and let $f$ be any function from length-$n$ inputs over
$\Omega$ to signs. The product resampling noise sensitivity of $f$ at resampling rate
zero is zero.
\end{theorem}

\begin{theorem}[Disagreement probability is at most one]
\lean{FABL.pmfDisagreementProbability_le_one}
\label{FABL.pmfDisagreementProbability_le_one}
\leanok
\uses{FABL.pmfExpectation, FABL.pmfExpectation_const, FABL.pmfExpectation_mono}
\difficulty{3}
Let $A$ be a finite set with a probability mass function, let $B$ be a set, and let
$f,g:A\to B$. The probability that $f$ and $g$ disagree on a sample drawn from this
distribution is at most $1$.
\end{theorem}

\begin{theorem}[Product resampling noise sensitivity le sqrt two]
\lean{FABL.productResamplingNoiseSensitivity_le_sqrt_two}
\statementsource{Analysis-of-Boolean-Functions-by-Ryan-ODonnell}{pdf-ffbe9c7b3306-p130-b004}
\label{FABL.productResamplingNoiseSensitivity_le_sqrt_two}
\leanok
\uses{FABL.IsProductLinearThreshold, FABL.Sign, FABL.pmfDisagreementProbability_le_one, FABL.productResamplingNoiseSensitivity, FABL.productResamplingNoiseSensitivity_le_peres}
\difficulty{5}
% dataset-pipeline: stub -- prose pending build_informal_statements.py
Statement pending regeneration.
\end{theorem}

\begin{theorem}[Product resampling noise sensitivity le sqrt two full]
\lean{FABL.productResamplingNoiseSensitivity_le_sqrt_two_full}
\statementsource{Analysis-of-Boolean-Functions-by-Ryan-ODonnell}{pdf-ffbe9c7b3306-p130-b004}
\label{FABL.productResamplingNoiseSensitivity_le_sqrt_two_full}
\leanok
\uses{FABL.IsProductLinearThreshold, FABL.Sign, FABL.productResamplingNoiseSensitivity, FABL.productResamplingNoiseSensitivity_le_sqrt_two, FABL.productResamplingNoiseSensitivity_zero}
\difficulty{4}
% dataset-pipeline: stub -- prose pending build_informal_statements.py
Statement pending regeneration.
\end{theorem}
